% GROK REVIEW::GROK-4.7-2222::2026-10-22T22:22:00+09:00
\documentclass[11pt]{article}
\usepackage[margin=1in]{geometry}
\usepackage{amsmath,amssymb}
\usepackage{url}
\title{Independent Grok 4.7 Review --- 2026-10-22 22:22 KST}
\author{Grok 4.7}
\date{22 October 2026}

\begin{document}
\maketitle

\begin{center}
\fbox{\parbox{0.9\linewidth}{
\textbf{Executive status: P versus NP remains OPEN.}
This review closes no load-bearing implication and changes no proof classification. It does not claim P$\neq$NP, P$=$NP, a Clay Millennium Prize result, or any unrestricted lower bound.
}}
\end{center}

\section{Scope and baseline}
Reviewed current HEAD of j15marti/pnp. Root \path{main.tex} is the minimal living document, previously updated at 21:19 KST by the prior independent review.

Also inspected:
\begin{itemize}
\item \path{AGENTS.md}, \path{CODEX.md}, \path{README.md};
\item structure and recent entries under \path{comments/grok-4.7/2026-10-22/} through the 21:19 review;
\item synthesis \path{research/main/2026-10-22_2148_MAIN_INTERFACE_ADAPTIVITY_BARRIER73.md}, which establishes that adaptive querying does not remove the exponential shared-interface factor in the detached black-box model (exact point-query evasiveness of $J\cdot 2^k$ and a prefix-existence lower bound of $2^{k-1}$ on a parity family);
\item synthesis \path{research/main/2026-10-22_2048_MAIN_INTERFACE_FLMV_RETUNING72.md}, which derives the exact FLMV feasibility surface for the restricted SMALL-INTERFACE64 constructor and a rational retuning point;
\item prior FLMV-PROBE-MARGIN71, PARENT-PROBE-BUDGET70, SHRUNK-QUERY-ENVELOPE67, SEMANTIC-PARENT-CALL-CIRCULARITY66, EQUALITY-LOCALIZATION-MINOR-BARRIER65, SOURCE-JOIN-COHERENCE56, and adversarial hash models;
\item confirmation that the protocol for independent reviews continues without scientific promotion. No load-bearing implication has been closed. External literature continues to show no accepted resolution of P versus NP.
\end{itemize}

\section{Observations}
\begin{enumerate}
\item Root \path{main.tex} remains a minimal, compilable status document with balanced ownership markers (LSPV-1 through LSPV-5), explicit OPEN status for P versus NP, the arbitrary-algorithm bridge, and unrestricted transfer of restricted ARC results.
\item INTERFACE-ADAPTIVITY-BARRIER73 carefully shows that adaptive ordering and prefix-existence queries do not asymptotically repair the detached small-interface compatibility problem. Point-membership access is fully evasive ($J\cdot 2^k$), and even blockwise prefix queries retain an exponential leaf lower bound on a simple parity family. The result is properly scoped as a detached-oracle barrier; source-aware escapes remain possible but must be explicitly charged. The constructive all-input hereditary transformer remains OPEN.
\item INTERFACE-FLMV-RETUNING72 provides a Pareto option around the original $(0.6,0.5)$ point, with a rational retuning permitting a slightly higher probe allowance at the cost of space and round exponents. It remains restricted accounting.
\item FLMV-PROBE-MARGIN71, PARENT-PROBE-BUDGET70, and prior barriers continue to provide disciplined restricted/computational results without promoting unrestricted claims.
\item No load-bearing implication from an arbitrary polynomial-time decider to a superpolynomial lower bound is closed. Relativizing and algebrizing barriers continue to apply. Restricted results do not transfer unrestricted. Progress now requires a charged source-aware mixed-child construction.
\item The comments/grok-4.7 folder structure is established and continues to receive independent timestamped reviews without altering scientific claims. The living document correctly references the most recent review.
\item External literature (as of this timestamp) shows no accepted proof resolving P versus NP; unverified claims continue to appear but do not affect the repository status.
\end{enumerate}

\section{Actions taken}
\begin{enumerate}
\item Added this review under \path{comments/grok-4.7/2026-10-22/}.
\item Updated root \path{main.tex} with a new review timestamp and a brief note confirming continued OPEN status and noting the adaptivity barrier without promoting any claim.
\end{enumerate}

\section{Recommendations}
\begin{enumerate}
\item Specialists should continue reconciling ownership regions against the total corpus, including INTERFACE-ADAPTIVITY-BARRIER73, INTERFACE-FLMV-RETUNING72, FLMV-PROBE-MARGIN71, PARENT-PROBE-BUDGET70, SHRUNK-QUERY-ENVELOPE67, SEMANTIC-PARENT-CALL-CIRCULARITY66, EQUALITY-LOCALIZATION-MINOR-BARRIER65, SMALL-INTERFACE64, SOURCE-JOIN-COHERENCE56, prior ARC restricted results, adversarial hash models, and any newer explorations under research/.
\item The highest-value open target remains construction of an exact, all-input, hereditary semantic quotient from strictly smaller charged queries of length $N^\beta$ for fixed $\beta<1$ (specifically the $(0.6,0.5,1.1,0.5)$ envelope, or the refined conditional $\eta\le 121/200$ slack, or the retuned Pareto point), surviving hostile controls, barrier audits (relativization, algebrization, natural proofs), model matching, exact quantifier accounting, aggregate parent-probe charges, and source-aware escapes from detached-interface barriers.
\item Maintain strict separation between finite computational evidence and asymptotic claims. Require published maximum query length, query count, aggregate parent probes, and source-acquisition costs for every proposed semantic compiler. Demand explicit source-aware modeling for any escape from detached-oracle lower bounds.
\end{enumerate}

\section{Status}
P versus NP remains open. No scientific claim was promoted.

% END_OF_VERIFIED_REVIEW::GROK-4.7::GROK-4.7-2222
\end{document}
